% R911: gerado; alterar dados e gerador, não esta prancha.
\clearpage
\refstepcounter{figure}\label{fig:AR16-structure}
\noindent{\large\sffamily\bfseries\color{navy}AR16 · Diagrama de arquitetura}\par\smallskip
\noindent\textbf{Infraestrutura, dados e segurança local}\par
\noindent\textit{Como executar verificações locais rastreáveis?}\par\medskip
\begin{center}

\tikzset{
 atlasnode/.style={draw=steel,fill=blueBg,rounded corners=2mm,align=center,
  text width=4cm,inner sep=5pt,font=\fontsize{9.5}{11.8}\selectfont,text=textmain},
 atlasroot/.style={atlasnode,fill=navy,text=white,text width=2.8cm,font=\fontsize{10}{12}\selectfont\bfseries},
 atlasdecision/.style={diamond,aspect=2.1,draw=amber,fill=amberBg,align=center,
  text width=3cm,inner sep=2pt,font=\fontsize{9.5}{11.5}\selectfont},
 atlascall/.style={-Stealth,draw=steel,thick},
 atlasevent/.style={-Stealth,draw=violet,dashed,thick},
 atlasdata/.style={-Stealth,draw=teal,densely dotted,thick},
 atlasstructure/.style={draw=slate,thick},
 atlasboundary/.style={-Stealth,draw=amber,dash dot,thick},
 atlasedgelabel/.style={fill=papel,inner sep=1.3pt,font=\fontsize{9}{10}\selectfont}
}

\begin{tikzpicture}
\node[atlasnode] (c0) at (-4.3,0) {\textbf{C1: Serviço de CI local}\\Aceita somente perfil científico fechado};
\node[atlasnode] (c1) at (4.3,0) {\textbf{C2: Snapshot de fontes}\\Identifica as fontes do perfil executado};
\node[atlasnode] (c2) at (-4.3,-3.0) {\textbf{C3: Supervisor de execução local}\\Observa runner e limites de recursos};
\node[atlasnode] (c3) at (4.3,-3.0) {\textbf{C4: Runner do perfil científico}\\Executa os grupos de testes enumerados};
\node[atlasnode] (c4) at (-4.3,-6.0) {\textbf{C5: Validador de recibos}\\Confere resultados dos grupos previstos};
\node[atlasnode] (c5) at (4.3,-6.0) {\textbf{C6: Consulta do estado local}\\Exibe status e compatibilidade de fontes};
\draw[atlasstructure] (c0.east) to[out=0,in=180] node[atlasedgelabel,pos=.25]{E1} (c1.west);
\draw[atlasboundary] (c0.south) -- node[atlasedgelabel,pos=.25]{E2} (c2.north);
\draw[atlasboundary] (c2.east) to[out=0,in=180] node[atlasedgelabel,pos=.25]{E3} (c3.west);
\draw[atlasstructure] (c3.west) to[out=180,in=0] node[atlasedgelabel,pos=.25]{E4} (c4.east);
\draw[atlasstructure] (c2.south) -- node[atlasedgelabel,pos=.25]{E5} (c4.north);
\draw[atlasdata] (c2.east) to[out=0,in=180] node[atlasedgelabel,pos=.25]{E6} (c5.west);
\end{tikzpicture}
\end{center}
{\small\textbf{Figura \thefigure.} Diagrama de arquitetura --- Infraestrutura, dados e segurança local}\par\smallskip
{\footnotesize\textbf{Chave das relações.} E1: C1 para C2 --- identifica fontes via workspace (estrutura); E2: C1 para C3 --- inicia supervisor em processo filho (fronteira); E3: C3 para C4 --- inicia runner do perfil fechado (fronteira); E4: C4 para C5 --- grupo confere os próprios recibos (estrutura); E5: C3 para C5 --- revalida via contrato de saída (estrutura); E6: C3 para C6 --- registro persistido de execução (dado).}\par\smallskip
\noindent\textbf{Como funciona e por quê.} Localize as seis responsabilidades C1--C6. Os identificadores E1--E6 permitem conferir o significado de cada conexão na chave das relações.\par
\noindent\textbf{Justificativa.} Separar componentes e tipos de relação ajuda a localizar onde uma dependência, um evento ou uma fronteira foi declarado.\par
\noindent\textbf{Limite.} A assinatura HMAC detecta alterações dos registros segundo o contrato local; não isola o proprietário do host. Essa CI não cria check no GitHub nem equivale a validação científica externa. Os testes cobrem o perfil enumerado. normalize, schemas\_validator e batch\_executor são utilitários simples, sem garantia implícita de segurança ou distribuição.\par\smallskip
\fontecore{integrations/local_ci.py}{LocalCIService.start}
\fontecore{integrations/local_ci.py}{LocalCIService._work}
